\chapter{Urticaria}
\index{Urticaria}

\section*{Definition and Overview}
Urticaria, historically derived from the Latin word ``urtica'' meaning stinging nettle, is a common and distressing dermatological reaction pattern characterized by the sudden development of transient wheals (hives), angioedema, or both. The hallmark of the condition is the presence of intensely pruritic, raised, erythematous, or pale-centered swelling of the skin. A defining clinical parameter of classic urticaria is its transient nature, where individual skin lesions develop rapidly, typically resolve completely within 24 hours without leaving any residual scarring or pigmentation, and may be replaced by new lesions in different anatomical locations. In approximately 40\% of affected individuals, urticaria co-occurs with angioedema, which involves a deeper, asymmetric swelling of the subcutaneous or submucosal tissues. Angioedema commonly affects the face, lips, tongue, extremities, or genitals, causing a painful, burning, or tight sensation rather than itching, and may take up to 72 hours to resolve.

Clinically, urticaria is classified based on its duration into acute and chronic forms. Acute urticaria refers to episodes that last for less than six weeks, often resolving within a few days to a couple of weeks. It is frequently triggered by acute infections (such as viral upper respiratory tract infections), immediate allergic reactions to foods, medications, or insect stings, or direct physical contact. Chronic urticaria is defined by the daily or near-daily presence of wheals for six weeks or longer. Chronic urticaria is further categorized into chronic spontaneous urticaria (CSU), where wheals appear without any identifiable external physical trigger, and chronic inducible urticaria (CIndU), where specific physical or environmental triggers, such as mechanical friction, cold, heat, pressure, solar radiation, or exercise, provoke the lesions. Due to the unpredictable nature of the condition and the intensity of the associated pruritus, urticaria imposes a significant physical and psychological burden on patients, impairing their sleep, daily activities, and overall quality of life to a degree comparable to that of severe cardiovascular or other debilitating chronic diseases.

\section*{Epidemiology}
Urticaria is a highly prevalent dermatological disorder with a substantial global public health footprint. Epidemiological investigations indicate that approximately 15\% to 25\% of the general population will experience at least one episode of acute urticaria during their lifetime. While acute urticaria can affect individuals of any age, it is particularly common in pediatric populations and young adults, often arising secondary to viral illnesses or sudden allergen exposures.

Chronic urticaria, although less common than the acute form, still affects approximately 0.5\% to 1.0\% of the global population at any given time. Chronic spontaneous urticaria demonstrates a significant gender disparity, affecting females approximately twice as frequently as males. The peak onset for chronic urticaria is observed between the ages of 30 and 50 years, which corresponds to the most economically productive years of life, thereby exacerbating the socioeconomic burden of the disease through direct medical costs and indirect productivity losses. Chronic inducible urticarias can affect all age groups, with symptomatic dermographism being the most common subtype, boasting an estimated prevalence of 2\% to 5\% in the general population. The persistence of chronic urticaria is a major clinical challenge; longitudinal studies show that about 50\% of patients remain symptomatic at one year, and 10\% to 20\% continue to experience recurrent flares for more than five years after their initial diagnosis.

\section*{Aetiology and Pathophysiology}
The development of urticaria involves a complex immunological and physiological cascade. The core mechanism centers on the activation and degranulation of mast cells in the skin, which release a variety of inflammatory mediators into the surrounding tissue.

\begin{itemize}
    \item \textbf{Mast Cell Activation and Mediator Release}: Mast cells, located in the perivascular space of the dermis, are the principal effector cells in urticaria. Upon activation, they undergo degranulation, releasing preformed mediators like histamine, heparin, and neutral proteases (tryptase and chymase), alongside newly synthesized lipid mediators (such as prostaglandin D2 and leukotrienes C4, D4, and E4) and cytokines (including tumor necrosis factor-alpha, interleukin-4, and interleukin-5).
    \item \textbf{Vascular Endothelial Changes}: Histamine binds to H1 and H2 receptors on vascular endothelial cells, triggering the release of nitric oxide and other vasodilatory signals. This leads to the rapid dilation of superficial dermal venules and a significant increase in capillary permeability. The subsequent extravasation of protein-rich plasma into the dermis creates the localized tissue edema that manifests clinically as a wheal.
    \item \textbf{Pruritus and Neurogenic Flare}: Histamine also binds to H1 receptors on unmyelinated sensory C-fibers in the skin, transmitting signals of intense itching to the central nervous system. The activation of these nerves induces an antidromic axon reflex, prompting the release of neuropeptides, such as substance P and calcitonin gene-related peptide (CGRP), which further promote local vasodilation and produce the erythematous flare surrounding the wheal.
    \item \textbf{Autoimmune Pathways (Type IIb and Type I Autoimmunity)}: In approximately 30\% to 40\% of cases of chronic spontaneous urticaria, the disease is driven by autoimmune mechanisms. Type IIb autoimmune urticaria involves circulating functional IgG autoantibodies directed against the alpha-subunit of the high-affinity IgE receptor (Fc-epsilon-RI) or against IgE itself. These autoantibodies cross-link the receptors on the surface of mast cells and basophils, triggering degranulation in the absence of external allergens. Type I autoimmune urticaria (autoallergic urticaria) involves IgE antibodies directed against autoantigens, such as thyroid peroxidase or double-stranded DNA.
    \item \textbf{IgE-Mediated Immediate Hypersensitivity}: In acute urticaria, type I immediate hypersensitivity reactions are common. Exposure to an exogenous allergen (such as penicillin, peanuts, shellfish, or bee venom) to which the patient has been previously sensitized causes the allergen to bind to and cross-link IgE antibodies already bound to Fc-epsilon-RI receptors on mast cells, leading to explosive degranulation.
    \item \textbf{Non-Immunological Mast Cell Activation}: Mast cells can also be activated directly through non-immunological pathways. Medications such as opioids, muscle relaxants, and radiocontrast media can trigger mast cell degranulation via specific receptors, such as the MRGPRX2 receptor. Furthermore, nonsteroidal anti-inflammatory drugs (NSAIDs) and aspirin can cause hives by inhibiting the cyclooxygenase-1 (COX-1) enzyme, which shifts arachidonic acid metabolism toward the production of pro-inflammatory leukotrienes.
    \item \textbf{Physical and Environmental Triggers}: In chronic inducible urticarias, mast cells exhibit hypersensitivity to physical forces. Mechanical friction (dermographism), exposure to cold air or water (cold urticaria), localized heat (heat urticaria), sustained pressure (delayed pressure urticaria), solar radiation (solar urticaria), or vibration (vibratory angioedema) act as direct triggers for mediator release.
\end{itemize}

\section*{Clinical Features}
The clinical presentation of urticaria is marked by distinct skin lesions and associated symptoms that can range from a few isolated spots to near-total body coverage.

\begin{itemize}
    \item \textbf{Wheals (Hives)}: The primary lesion is the wheal, which is a raised, swollen, well-demarcated area of skin. Wheals are characteristically pale or white in the center due to localized edema compressing dermal blood vessels, surrounded by a bright red, erythematous flare. They vary in size from small papules of a few millimeters to large, geographic plaques that can migrate and merge.
    \item \textbf{Pruritus (Itching)}: Intense itching is the most prominent and distressing symptom of urticaria. The sensation is often described as an intense itch, sting, or burn. It typically worsens during the evening and night, or when the skin is warmed by blankets, hot baths, physical exertion, or emotional stress, leading to significant sleep disruption.
    \item \textbf{Lesion Transience}: A key clinical feature is that individual wheals are highly transient, appearing rapidly and resolving completely within 24 hours. The skin returns to its normal appearance without any residual scaling, bruising, or pigmentary changes, even though new wheals may continue to develop in other areas.
    \item \textbf{Dermographism}: In many individuals, particularly those with physical urticaria, lightly scratching or stroking the skin produces a linear wheal and flare along the line of friction within minutes. This is commonly referred to as ``skin writing''.
    \item \textbf{Angioedema}: Swelling of the deeper subcutaneous or submucosal tissues occurs in association with urticaria in approximately 40\% of cases. It is characterized by poorly defined, asymmetric swelling that is often painful, tight, or burning rather than itchy. It commonly affects the lips, eyelids, tongue, hands, feet, and genitalia, and can take up to 72 hours to resolve.
    \item \textbf{Systemic Symptoms}: In severe acute cases, or when hives are part of a broader systemic reaction, patients may experience mild systemic signs such as low-grade fever, headache, joint pain (arthralgia), abdominal cramping, nausea, vomiting, or diarrhea.
\end{itemize}

\section*{Differential Diagnosis}
Because many skin conditions can present with red, itchy rashes, it is essential to distinguish urticaria from other dermatological and systemic diseases to ensure correct treatment.

\begin{itemize}
    \item \textbf{Urticarial Vasculitis}: Unlike classic urticaria, the lesions of urticarial vasculitis remain fixed in the same anatomical location for more than 24 hours. They are typically painful, burning, or tender rather than itchy, and they often resolve leaving residual purpura (bruising), hyperpigmentation, or fine scaling. A skin biopsy is diagnostic, showing leukocytoclastic vasculitis.
    \item \textbf{Erythema Multiforme}: This acute inflammatory condition presents with characteristic target or ``bullseye'' lesions. These lesions have a dusky, blistered, or necrotic center, surrounded by a pale ring and an erythematous outer border. They are fixed for several days, are often symmetrical, and typically involve the palms, soles, and mucous membranes.
    \item \textbf{Atopic Dermatitis (Eczema)}: Atopic dermatitis is a chronic, relapsing skin condition characterized by dry, scaly, and poorly demarcated erythematous plaques. Unlike the transient, smooth wheals of urticaria, eczematous lesions are persistent, show signs of chronic scratching (excoriations and lichenification), and typically affect the flexural creases (elbows and knees).
    \item \textbf{Contact Dermatitis}: This is a localized, fixed inflammatory skin reaction caused by direct contact with an allergen (allergic contact dermatitis) or an irritant (irritant contact dermatitis). The resulting rash is limited to the site of contact, characterized by erythema, papules, vesicles, scaling, and itching, and persists for days to weeks.
    \item \textbf{Bullous Pemphigoid (Prodromal Stage)}: In the early prodromal stage of this autoimmune blistering disease, patients can present with intensely pruritic, urticarial-like plaques. However, these plaques are fixed and eventually develop tense, fluid-filled blisters (bullae) on top of or adjacent to them.
    \item \textbf{Mastocytosis (Urticaria Pigmentosa)}: This condition is characterized by abnormal accumulations of mast cells in the skin, appearing as fixed, reddish-brown macules and papules. When rubbed, the lesions swell and become itchy (Darier's sign), but the underlying brown spots remain permanently visible and do not resolve.
    \item \textbf{Anaphylaxis}: Although urticaria and angioedema are common features of anaphylaxis, anaphylaxis is a life-threatening, multi-systemic allergic reaction. It is differentiated by the rapid co-occurrence of respiratory distress, cardiovascular instability (hypotension, tachycardia, shock), or severe gastrointestinal symptoms.
\end{itemize}

\section*{When to Seek Medical Care}
Although most episodes of urticaria are benign, self-limiting, and easily managed, the condition can occasionally be a component of a severe, life-threatening systemic allergic reaction. Immediate medical evaluation is required under specific circumstances.

Patients must seek emergency medical services immediately (such as calling local emergency services in many countries, 911 in the United States, or 999 in the United Kingdom) if they experience hives or swelling accompanied by any of the following symptoms:
\begin{itemize}
    \item Difficulty breathing, shortness of breath, audible wheezing, or a high-pitched whistling sound when breathing (stridor).
    \item Swelling of the tongue, throat, or mouth that interferes with breathing or makes it difficult to swallow.
    \item Sudden hoarseness, a changing voice, or a feeling of tightness in the throat.
    \item Dizziness, lightheadedness, confusion, feeling faint, or loss of consciousness, which indicate a severe drop in blood pressure.
    \item Severe abdominal pain, nausea, vomiting, or diarrhea occurring rapidly alongside the skin rash.
\end{itemize}

Non-emergency medical care from a general practitioner or a dermatologist should be sought if:
\begin{itemize}
    \item The hives do not improve within 48 hours of onset despite taking over-the-counter, non-sedating antihistamines.
    \item The symptoms are severe enough to interfere significantly with sleep, work, or daily activities.
    \item Individual lesions persist in the same exact spot for more than 24 hours, or resolve leaving behind bruising, discoloration, or scaling.
    \item The hives continue to appear regularly for more than six weeks (suggesting chronic urticaria).
    \item The rash is accompanied by joint pain, high fever, or general malaise.
\end{itemize}

\section*{Diagnosis and Investigations}
The diagnosis of urticaria is primarily clinical, based on a comprehensive medical history and a thorough physical examination. Extensive diagnostic testing is rarely necessary for acute urticaria unless a specific allergen is strongly suspected. For chronic urticaria, investigations are focused on ruling out underlying systemic diseases and confirming specific inducible types.

\begin{itemize}
    \item \textbf{Clinical History}: A detailed history is the most important diagnostic tool. The clinician will investigate the frequency and duration of lesions, the presence of itching versus pain, any associated angioedema, potential triggers (foods, medications, physical stimuli, infections, insect stings), travel history, and a family history of allergic or autoimmune conditions.
    \item \textbf{Physical Examination}: The doctor will inspect the skin to confirm the presence of classic wheals. A physical exam may include physical provocation testing, such as stroking the skin with a blunt instrument to check for dermographism, applying an ice cube to the forearm to test for cold urticaria, or applying localized pressure to check for delayed pressure urticaria.
    \item \textbf{Basic Laboratory Screening}: In cases of chronic spontaneous urticaria, a limited set of blood tests is recommended to rule out underlying systemic disorders. This typically includes a Full Blood Count (FBC) to check for eosinophilia or leukocytosis, and inflammatory markers such as C-Reactive Protein (CRP) or Erythrocyte Sedimentation Rate (ESR).
    \item \textbf{Allergy Testing}: Skin prick testing or serum allergen-specific IgE blood tests may be performed if the history strongly suggests a specific type I IgE-mediated allergy (e.g., to a particular food or environmental allergen). Routine, indiscriminate allergy screening is not recommended.
    \item \textbf{Thyroid Evaluation}: Because chronic spontaneous urticaria is associated with autoimmune thyroid disease, testing for thyroid-stimulating hormone (TSH) and thyroid autoantibodies (anti-thyroid peroxidase and anti-thyroglobulin) can help identify this association.
    \item \textbf{Autologous Serum Skin Test (ASST)}: This in vivo test is sometimes used to screen for autoantibodies against IgE or the IgE receptor. A small sample of the patient's own serum is isolated and injected intradermally; a subsequent wheal and flare reaction indicates the presence of circulating functional autoantibodies.
    \item \textbf{Skin Biopsy}: A punch biopsy of a representative skin lesion is indicated if the hives are atypical. Biopsy is essential if lesions last more than 24 hours, are painful rather than itchy, or leave residual purpura, to rule out urticarial vasculitis.
\end{itemize}

\section*{Management}
The management of urticaria focuses on symptom control and, when possible, the identification and avoidance of known triggers. A step-care approach is recommended by international guidelines, focusing on pharmacological therapies that suppress mast cell mediator release and action.

\begin{itemize}
    \item \textbf{First-Line Pharmacotherapy (Standard-Dose H1 Antihistamines)}: The cornerstone of urticaria treatment is the daily use of non-sedating, second-generation H1 antihistamines (e.g., cetirizine, loratadine, fexofenadine, desloratadine, or bilastine). These medications effectively block H1 receptors on blood vessels and nerves, reducing wheals and pruritus. First-generation antihistamines (such as diphenhydramine or promethazine) are generally avoided due to their significant sedative and anticholinergic side effects.
    \item \textbf{Second-Line Pharmacotherapy (Up-dosing of H1 Antihistamines)}: If standard doses of second-generation antihistamines do not control symptoms within 2 to 4 weeks, the dose can be increased up to four times the standard licensed dose (e.g., fexofenadine 720 mg daily or cetirizine 40 mg daily) under medical supervision. This approach is highly effective and remains safe for most patients.
    \item \textbf{Third-Line Pharmacotherapy (Add-on Omalizumab)}: For patients with chronic urticaria who remain refractory to high-dose antihistamines, the monoclonal antibody omalizumab is the gold standard add-on therapy. Omalizumab binds to free IgE, preventing it from binding to mast cell receptors and downregulating receptor expression. It is administered as a subcutaneous injection, typically every 4 weeks.
    \item \textbf{Fourth-Line Pharmacotherapy (Immunosuppressants)}: In severe, refractory cases of chronic urticaria that fail to respond to omalizumab, oral immunosuppressive agents may be considered. Cyclosporine A is the most widely supported agent in this category, as it inhibits mast cell degranulation. However, its use is limited by potential side effects, including renal toxicity and hypertension, requiring close monitoring.
    \item \textbf{Short-Term Systemic Corticosteroids}: A brief course of oral corticosteroids (such as prednisolone) may be prescribed to control severe exacerbations of acute or chronic urticaria. Their use should be strictly limited to a maximum of 3 to 10 days to avoid long-term systemic side effects, such as adrenal suppression, osteoporosis, and weight gain.
    \item \textbf{Emergency Epinephrine}: For acute urticaria presenting as part of a systemic anaphylactic reaction, intramuscular epinephrine is the immediate, life-saving treatment of choice. Patients with a known risk of anaphylaxis must carry an auto-injector (e.g., EpiPen) at all times.
\end{itemize}

\section*{Complications}
Although urticaria itself is rarely fatal, it can lead to several physical, psychological, and systemic complications, particularly when it becomes chronic or is accompanied by angioedema.

\begin{itemize}
    \item \textbf{Airway Compromise}: The most severe acute complication is respiratory distress caused by laryngeal angioedema. Swelling of the tissues in the throat and larynx can rapidly obstruct the airway, leading to asphyxiation if not treated immediately with epinephrine.
    \item \textbf{Psychological Distress and Sleep Deprivation}: Chronic urticaria has a profound impact on mental health. The unpredictable nature of the flares, combined with relentless nocturnal itching, leads to severe sleep deprivation. This, in turn, increases the risk of developing anxiety, depression, chronic fatigue, and diminished work or academic performance.
    \item \textbf{Secondary Skin Infections}: Intense scratching due to severe pruritus can compromise the skin barrier. This leads to excoriations (scratches), which can become infected by bacteria such as \textit{Staphylococcus aureus}, resulting in impetigo, cellulitis, or localized abscesses.
    \item \textbf{Anaphylaxis and Shock}: In cases where acute urticaria is a manifestation of systemic anaphylaxis, the rapid and widespread release of mediators can cause profound vasodilation, leading to cardiovascular collapse, severe hypotension, and anaphylactic shock.
    \item \textbf{Side Effects of Long-Term Medication}: While antihistamines are relatively safe, chronic use of high doses can cause mild drowsiness or dry mouth. More significantly, inappropriate long-term use of systemic corticosteroids to manage chronic hives can lead to severe complications, including Cushingoid features, diabetes mellitus, hypertension, osteoporosis, cataracts, and increased susceptibility to infections.
\end{itemize}

\section*{Prognosis and Outcomes}
The prognosis for urticaria is generally excellent, particularly for the acute form, which is typically transient and resolves spontaneously.

In patients with acute urticaria, the vast majority of episodes resolve completely within a few days to a couple of weeks. Once the triggering allergen, infection, or physical stimulus has cleared, the skin returns to normal without any permanent damage or scarring. Only a small fraction of acute cases transition into chronic urticaria.

For chronic urticaria, the prognosis is more variable but remains favorable overall. While the condition can persist for years, it is characterized by high rates of spontaneous remission. Studies indicate that approximately 30\% to 50\% of patients with chronic spontaneous urticaria experience remission within one year of onset. An additional 20\% to 30\% achieve remission within three to five years. Factors that correlate with a longer duration of chronic urticaria include severe initial symptoms, the co-occurrence of angioedema, positive autologous serum skin tests, and the presence of physical or inducible urticarias. Fortunately, even for patients with long-lasting disease, modern therapeutic options like omalizumab and high-dose second-generation antihistamines allow the vast majority to achieve complete or near-complete control of their symptoms, maintaining a normal quality of life until the disease naturally resolves.

\section*{Prevention and Self-Care}
While it is not always possible to prevent urticaria, especially when the cause is spontaneous or related to an underlying autoimmune process, several self-care measures and prevention strategies can help manage symptoms and reduce the frequency of flare-ups.

\begin{itemize}
    \item \textbf{Symptom Diary Maintenance}: Patients should keep a detailed daily journal recording food intake, medication use, physical activities, environmental exposures, stress levels, and the severity of hives to help identify and eliminate potential triggers.
    \item \textbf{Avoidance of Common Exacerbating Substances}: Patients with chronic urticaria should avoid aspirin and other nonsteroidal anti-inflammatory drugs (NSAIDs) such as ibuprofen, as these can provoke flares. Acetaminophen (paracetamol) is a safer alternative for pain relief and fever.
    \item \textbf{Cooling and Soothing Techniques}: Applying cool, damp compresses or taking cool baths can help constrict dermal capillaries, decrease swelling, and relieve intense itching. Hot water should be avoided, as it promotes vasodilation and worsens pruritus.
    \item \textbf{Skin Barrier Support}: Daily application of mild, fragrance-free, hypoallergenic moisturizers helps protect the skin barrier and reduces overall skin sensitivity. Calamine lotion can be applied to individual itchy wheals for soothing relief.
    \item \textbf{Friction and Pressure Reduction}: To prevent hives triggered by mechanical forces, patients should wear loose-fitting, soft clothing made of natural fibers like cotton, and avoid tight belts, heavy backpack straps, or rough synthetic fabrics.
    \item \textbf{Stress Reduction Practices}: Because stress is a powerful trigger for mast cell degranulation, practicing stress-reduction techniques such as mindfulness, meditation, progressive muscle relaxation, deep breathing exercises, or seeking psychological support can help lower the frequency of flares.
    \item \textbf{Alcohol and Spicy Food Limitation}: Alcohol and highly spiced foods can cause cutaneous vasodilation and exacerbate itching, and their consumption should be minimized during active flare-ups.
\end{itemize}

\section*{Sources and Further Reading}
\begin{itemize}
    \item Australasian Society of Clinical Immunology and Allergy (ASCIA). Urticaria (Hives) Patient Information Guide. Available at: https://www.allergy.org.au
    \item World Allergy Organization (WAO). Global Consensus Guidelines for the Management of Urticaria. Available at: https://www.worldallergy.org
    \item DermNet New Zealand. Comprehensive Guide on Urticaria (Hives). Available at: https://www.dermnetnz.org
    \item American Academy of Dermatology (AAD). Hives: Overview, Causes, and Treatment. Available at: https://www.aad.org
    \item Mayo Clinic. Chronic Hives (Urticaria) - Symptoms and Causes. Available at: https://www.mayoclinic.org
    \item National Health Service (NHS). Hives (Urticaria) Overview and Management. Available at: https://www.nhs.uk
\end{itemize}